  % Versão resumida do relatório. Mantenha resultados.tex no mesmo diretório
% para que as macros numéricas sejam carregadas.
\documentclass[11pt,a4paper]{article}

\usepackage[utf8]{inputenc}
\usepackage[T1]{fontenc}
\usepackage[brazil]{babel}
\usepackage{lmodern}
\usepackage{geometry}
\geometry{margin=2.2cm}
\usepackage{amsmath}
\usepackage{booktabs,tabularx}
\usepackage{hyperref}
\hypersetup{colorlinks=true,linkcolor=black,citecolor=black,urlcolor=blue}

\newcommand{\Qdot}{\dot Q}
\newcommand{\Patm}{917{,}2}
\newcommand{\Tsat}{97{,}2}
\input{resultados.tex}

\begin{document}

\begin{center}
{\small UNIVERSIDADE FEDERAL DE MINAS GERAIS\\
Escola de Engenharia -- Engenharia Aeroespacial\\
EMA-103 -- Laboratório de Térmica}

\vspace{0.8cm}
{\Large\bfseries Transferência de calor em uma aleta cilíndrica}

\vspace{0.6cm}
\begin{tabular}{rl}
\textbf{Aluno(a):} & Rodrigo Malato Navegantes\\
\textbf{Matrícula:} & 2023038175\\
\textbf{Turma:} & EA1\\
\textbf{Professor(a):} & Luiz Machado\\
\textbf{Data da prática:} & 13 de setembro de 2026
\end{tabular}
\end{center}

\section{Objetivo}

Avaliar o perfil de temperatura e a transferência de calor de uma aleta
cilíndrica de aço aquecida por uma câmara de vapor. Os resultados medidos são
comparados ao modelo de aleta reta de seção uniforme.

\section{Modelo utilizado}

Admitiu-se regime permanente, condução predominantemente axial, propriedades
constantes e coeficiente médio de troca térmica $h$. A temperatura varia ao
longo do comprimento da aleta, mas é aproximadamente uniforme em cada seção
transversal. Para uma aleta cilíndrica de diâmetro $d$, comprimento $L$ e
condutividade $k$, o perímetro e a área transversal são
\begin{equation}
  P=\pi d,\qquad A=\frac{\pi d^2}{4},
\end{equation}
e o parâmetro que relaciona a condução na aleta com as perdas para o ambiente é
\begin{equation}
  m=\sqrt{\frac{hP}{kA}}=\sqrt{\frac{4h}{kd}}.
\end{equation}

Definindo $\theta=T-T_{\mathrm{amb}}$, o balanço de energia em um elemento da
aleta resulta em
\begin{equation}
  \frac{d^2\theta}{dx^2}-m^2\theta=0.
\end{equation}
Na base, a condição é $T(0)=T_b$. Na ponta, foi considerada convecção pela área
da seção transversal. Com essas condições, o perfil previsto é~\cite{incropera}
\begin{equation}
\frac{T(x)-T_{\mathrm{amb}}}{T_b-T_{\mathrm{amb}}}=
\frac{\cosh[m(L-x)]+\dfrac{h}{mk}\sinh[m(L-x)]}
{\cosh(mL)+\dfrac{h}{mk}\sinh(mL)}.
\label{eq:perfil}
\end{equation}

A taxa teórica foi obtida pela condução na raiz da aleta. Como conferência, a
taxa experimental foi calculada integrando as perdas convectivas e radiativas
representadas por $h$ ao longo do perfil de temperaturas medido. A efetividade $\epsilon$
compara a aleta com a base sem aleta; a eficiência $\eta$ compara a taxa real com a
taxa que ocorreria se toda a aleta estivesse na temperatura da base.

\section{Procedimento e dados}

Foram usados uma aleta cilíndrica de aço, câmara e gerador de vapor, termopar
tipo K, pirômetro, escala e paquímetro. Após a estabilização do conjunto,
mediram-se a temperatura da base, a temperatura ambiente e as temperaturas ao
longo da aleta nas posições de 45 a 295 mm. O comprimento e o diâmetro foram
medidos diretamente; foram adotados ainda $k=(60\pm5)$ W/(m K) e
$h=(12\pm2)$ W/(m$^2$ K), conforme o roteiro da disciplina~\cite{roteiro}.

Como a câmara é alimentada por vapor, adotou-se a pressão atmosférica registrada
às 17h (horário de Brasília) pela estação meteorológica
\emph{Belo Horizonte--Pampulha} do INMET: \Patm\ hPa~\cite{inmet}. Essa pressão
corresponde a uma temperatura de saturação da água de aproximadamente
\Tsat\ $^\circ$C.

Os cálculos, a propagação das incertezas e as tabelas foram gerados pelo arquivo
\href{}{\texttt{codigo-03.py}} do repositório da disciplina~\cite{codigo}.

\begin{table}[h]
\centering
\caption{Dados usados na análise.}
\small
\begin{tabular}{lcc}
\toprule
\textbf{Grandeza} & \textbf{Valor} & \textbf{Incerteza ou observação}\\
\midrule
$L$ (mm) & \Laleta & $\pm0{,}5$ mm\\
$d$ (mm) & \Daleta & $\pm0{,}1$ mm\\
$T_b=T(0)$ ($^{\circ}$C) & \Tzero & $\pm1{,}0\,^{\circ}$C\\
$T(45)$ ($^{\circ}$C) & \Tquarentacinco & $\pm1{,}0\,^{\circ}$C\\
$T(95)$ ($^{\circ}$C) & \Tnoventacinco & $\pm1{,}0\,^{\circ}$C\\
$T(145)$ ($^{\circ}$C) & \Tcentoquarentacinco & $\pm1{,}0\,^{\circ}$C\\
$T(195)$ ($^{\circ}$C) & \Tcentonoventacinco & $\pm1{,}0\,^{\circ}$C\\
$T(245)$ ($^{\circ}$C) & \Tduzentosquarentacinco & $\pm1{,}0\,^{\circ}$C\\
$T(295)$ ($^{\circ}$C) & \Tduzentosnoventacinco & $\pm1{,}0\,^{\circ}$C\\
$T_{\mathrm{amb}}$ ($^{\circ}$C) & \Tamb & $\pm1{,}0\,^{\circ}$C\\
$p_{\mathrm{atm}}$ (hPa) & \Patm & ---\\
\bottomrule
\end{tabular}
\label{tab:medicoes}
\end{table}

As incertezas foram propagadas numericamente a partir das incertezas dos dados
de entrada. Os resultados são apresentados com incerteza expandida ($k=2$).

\section{Resultados e discussão}

\subsection{Perfil de temperaturas}

A Tabela~\ref{tab:perfil} compara as temperaturas medidas com a
Equação~\eqref{eq:perfil}. O ponto na base é condição de contorno do modelo;
os demais pontos são usados para avaliar a concordância.

\begin{table}[h]
\centering
\caption{Comparação entre os perfis experimental e teórico.}
\small
\begin{tabular}{ccccc}
\toprule
$x$ (mm) & $T_{\mathrm{exp}}$ ($^{\circ}$C) &
$T_{\mathrm{teo}}\pm U$ ($^{\circ}$C) &
$\delta_T$ ($^{\circ}$C) & Compatível?\\
\midrule
0   & \Tzero & \Tteozero$\pm$\UTteozero & \Deltazero & \Compzero\\
45  & \Tquarentacinco & \Tteoquarentacinco$\pm$\UTteoquarentacinco & \Deltaquarentacinco & \Compquarentacinco\\
95  & \Tnoventacinco & \Tteonoventacinco$\pm$\UTteonoventacinco & \Deltanoventacinco & \Compnoventacinco\\
145 & \Tcentoquarentacinco & \Tteocentoquarentacinco$\pm$\UTteocentoquarentacinco & \Deltacentoquarentacinco & \Compcentoquarentacinco\\
195 & \Tcentonoventacinco & \Tteocentonoventacinco$\pm$\UTteocentonoventacinco & \Deltacentonoventacinco & \Compcentonoventacinco\\
245 & \Tduzentosquarentacinco & \Tteoduzentosquarentacinco$\pm$\UTteoduzentosquarentacinco & \Deltaduzentosquarentacinco & \Compduzentosquarentacinco\\
295 & \Tduzentosnoventacinco & \Tteoduzentosnoventacinco$\pm$\UTteoduzentosnoventacinco & \Deltaduzentosnoventacinco & \Compduzentosnoventacinco\\
\bottomrule
\end{tabular}
\label{tab:perfil}
\end{table}

O perfil apresentou concordância \QualificacaoPerfil. Dos seis pontos
independentes, \NCompativeis{} são compatíveis dentro das incertezas e
\NRoteiro{} atendem ao critério de diferença inferior a $1\,^{\circ}$C. O maior
desvio ocorreu em $x=\XMaiorDesvio$ mm e valeu
\MaiorDesvio\,$^{\circ}$C. Como os demais pontos acompanham a tendência prevista,
esse resultado sugere uma perturbação localizada de leitura, como contato
imperfeito do termopar ou corrente de ar.

\subsection{Transferência de calor e desempenho}

\begin{table}[h]
\centering
\caption{Resultados principais da aleta. Incertezas expandidas, com $k=2$.}
\small
\begin{tabular}{lcc}
\toprule
\textbf{Grandeza} & \textbf{Resultado} & \textbf{Unidade}\\
\midrule
Produto $mL$ & \mL & --\\
Taxa teórica, $\Qdot_{\mathrm{teo}}$ & $\Qteorico\pm\UQteorico$ & W\\
Taxa pelo perfil medido, $\Qdot_{\mathrm{exp}}$ & $\Qexperimental\pm\UQexperimental$ & W\\
Desvio entre as taxas & \DesvioQ & \%\\
Efetividade teórica, $\varepsilon_f$ & $\EfetividadeTeorica\pm\UEfetividadeTeorica$ & --\\
Eficiência teórica, $\eta_f$ & $\EficienciaTeorica\pm\UEficienciaTeorica$ & --\\
\bottomrule
\end{tabular}
\label{tab:desempenho}
\end{table}

As taxas de calor foram \QualificacaoTaxa: a estimativa teórica foi
$\Qteorico\pm\UQteorico$ W e a estimativa obtida do perfil medido foi
$\Qexperimental\pm\UQexperimental$ W. A diferença foi de \DesvioQ\%. Como as
duas estimativas usam os mesmos valores de $h$ e de geometria, essa comparação
mostra consistência entre o perfil e o modelo, mas não é uma validação totalmente
independente da taxa de calor.

A eficiência de $\EficienciaTeorica\pm\UEficienciaTeorica$ é limitada pela queda
de temperatura ao longo da aleta. Ainda assim, a efetividade de
$\EfetividadeTeorica\pm\UEfetividadeTeorica$ mostra que a aleta transfere muito
mais calor do que a base sozinha. Portanto, uma eficiência moderada não impede
que a aleta seja útil.

\section{Conclusão}

O modelo de aleta cilíndrica representou adequadamente as temperaturas medidas:
\NCompativeis{} dos seis pontos independentes foram compatíveis dentro das
incertezas. As taxas de calor teórica e experimental também foram
\QualificacaoTaxa, com desvio de \DesvioQ\%.

A pressão atmosférica adotada foi a leitura das 17h da estação Pampulha do
INMET: \Patm\ hPa. Esse valor implica temperatura de saturação de aproximadamente
\Tsat\,$^{\circ}$C e serve como verificação da fonte de vapor. A principal
limitação da análise continua sendo a incerteza do coeficiente de troca térmica
$h$.

\begin{thebibliography}{9}

\bibitem{roteiro}
UFMG. \textit{EMA-103: Laboratório de Térmica -- Aulas 6 e 7: aleta
cilíndrica}. Material didático da disciplina, 2026.

\bibitem{incropera}
BERGMAN, T. L.; LAVINE, A. S.; INCROPERA, F. P.; DEWITT, D. P.
\textit{Fundamentals of Heat and Mass Transfer}. 7. ed. Hoboken: Wiley, 2011.

\bibitem{inmet}
INSTITUTO NACIONAL DE METEOROLOGIA (INMET). \textit{Dados horários da estação
automática Belo Horizonte--Pampulha, 13 set. 2026}. Plataforma WIS2, consulta
em 21 set. 2026. \url{https://wis2bra.inmet.gov.br/oapi/collections}.

\bibitem{codigo}
NAVEGANTES, R. M. \textit{LabTermo: código da Prática 3}. Arquivo
\texttt{pratica-03/codigo-03.py}. Disponível em:
\url{https://github.com/rodrigo-navegantes/labtermo}.

\end{thebibliography}

\end{document}
